\chapter{Implementazione passo-passo}

Questo capitolo descrive il flusso operativo per ogni funzionalita, con frammenti di codice e spiegazioni dettagliate.

\section{Inizializzazione globale}
Nel file principale i moduli sono istanziati con parametri chiave e con i pin definiti in modo centralizzato. Questa scelta rende immediata la lettura della configurazione.

\begin{lstlisting}[style=cppstyle]
const int PIN_SERVOMOTOR = 3;
const int PIN_GREEN_LIGHT = 4;
const int PIN_RED_LIGHT = 5;

Turnstile entranceTurnstile(PIN_IN_BUTTON, PIN_OUT_BUTTON, 5);
Stoplight trafficLight(PIN_GREEN_LIGHT, PIN_RED_LIGHT);
Thermostat mainThermostat(PIN_TEMPERATURE_SENSOR, 20.0, PIN_YELLOW_LIGHT, PIN_BLU_LIGHT);
LightingControl galleryLighting(PIN_PHOTORESISTOR, PIN_PLAFONIERE, 200);
HumidifierControl galleryHumidifier(IS_HUMIDIFIER_LED, 65.0);
DisplayPanel galleryDisplay(PIN_LCD_RS, PIN_LCD_EN, PIN_LCD_D4, PIN_LCD_D5, PIN_LCD_D6, PIN_LCD_D7, 5);
\end{lstlisting}

\paragraph{Motivazione} I target e la capienza sono centralizzati per ridurre errori e per consentire una taratura rapida.

\section{Tornello e conteggio persone}
Il tornello legge i pulsanti e muove il servo con protezione meccanica.

\begin{lstlisting}[style=cppstyle]
if (inState == HIGH && _currentPeople < _maxPeople) {
    moveServo(90);
    _currentPeople++;
    delay(1000);
    moveServo(-90);
}
\end{lstlisting}

\paragraph{Passi} 1) lettura ingressi, 2) controllo capienza, 3) apertura servo, 4) aggiornamento contatore, 5) chiusura servo.

\paragraph{Sicurezza} Il contatore evita ingressi oltre il limite e la funzione \texttt{antiSufferingServo} impedisce movimenti fuori range.

\section{Semaforo di capienza}
Il semaforo traduce il conteggio in uno stato immediatamente leggibile.

\begin{lstlisting}[style=cppstyle]
if (currentPeople < maxPeople) {
    led(_greenPin, HIGH);
    led(_redPin, LOW);
} else {
    led(_greenPin, LOW);
    led(_redPin, HIGH);
}
\end{lstlisting}

\paragraph{Motivazione} Due soli stati riducono l'ambiguita per l'operatore.

\section{Termostato}
Il termostato legge la temperatura e decide l'azione con tolleranza e pausa anti-toggle.

\begin{lstlisting}[style=cppstyle]
if (temperature_celsius < (_desiredTemperature - tolerance)) {
    desired_action = 1;
} else if (temperature_celsius > (_desiredTemperature + tolerance)) {
    desired_action = 2;
} else {
    desired_action = 0;
}
\end{lstlisting}

\paragraph{Passi} 1) lettura sensore, 2) validazione, 3) confronto con target, 4) scelta azione, 5) attuazione con pausa.

\paragraph{Sicurezza} In caso di lettura invalida, entrambe le uscite sono spente e il metodo ritorna \texttt{false}.

\section{Controllo umidita}
Il modulo confronta la lettura del DHT22 con il target e una tolleranza fissa.

\begin{lstlisting}[style=cppstyle]
if (currentHumidity > (_targetHumidity + tolerance)) {
    led(_humidifierPin, HIGH);
} else if (currentHumidity < (_targetHumidity - tolerance)) {
    led(_humidifierPin, LOW);
} else {
    led(_humidifierPin, LOW);
}
\end{lstlisting}

\paragraph{Sicurezza} Se il sensore ritorna \texttt{NAN}, la lettura diventa \texttt{-999.0} e il modulo spegne l'uscita.

\section{Illuminazione}
Il controllo luci usa un modello semplice per tradurre la lettura analogica in lux e decide il PWM.

\begin{lstlisting}[style=cppstyle]
float missingLight = _targetLux - currentLux;
if (missingLight > 0) {
    pwm_value = 255;
} else {
    pwm_value = 0;
}
ledDimming(_dimmerPin, pwm_value);
\end{lstlisting}

\paragraph{Motivazione} In simulazione Wokwi si usa un PWM pieno quando serve luce, per rendere evidente la risposta.

\section{Display LCD}
Il display ruota pagine con informazioni diverse e mostra gli errori con priorita.

\begin{lstlisting}[style=cppstyle]
if (now - _lastPageChange >= 5000) {
    _page = (_page + 1) % 3;
    _lastPageChange = now;
}
\end{lstlisting}

\paragraph{Passi} 1) limitazione refresh, 2) rotazione pagina, 3) lettura valori dai moduli, 4) stampa pagina selezionata.

\paragraph{Sicurezza} Le letture invalide sono mostrate come \texttt{ERR} o come messaggi di errore, evitando interpretazioni errate.
